\documentclass[11pt,a4paper]{article}

\usepackage[margin=1in]{geometry}
\usepackage{hyperref}
\usepackage{enumitem}
\usepackage{booktabs}
\usepackage{tabularx}
\usepackage{xcolor}
\usepackage{mdframed}

\hypersetup{colorlinks=true, linkcolor=blue!70!black, urlcolor=blue!70!black}

\title{Current Codebase Deep Read\\[0.3em]
\large Read-only summary of the current SDK, runtime, and MBRA paths}
\author{Project Working Draft}
\date{\today}

\begin{document}

\maketitle

\section{Purpose}

This note is a \textbf{read-only understanding document} for the current state
of the repository after many iterative changes.

The goal is not to propose a redesign. The goal is to capture what the code
\textbf{actually does now}, especially in the following areas:

\begin{itemize}[nosep]
  \item the SDK bridge,
  \item the live runtime,
  \item the current simple-controller path,
  \item the localization / temporal state path,
  \item the current MBRA deployment wrapper,
  \item the optional depth-safety path.
\end{itemize}


\section{High-Level Bottom Line}

\begin{mdframed}
The repo is no longer just ``baseline localization + graph planner + a small
controller.'' It is now a more integrated system with:

\begin{itemize}[nosep]
  \item a \textbf{hybrid SDK bridge} that supports both browser-relay and direct RTM fallback,
  \item a much heavier \textbf{live runtime wrapper} with recovery and guardrails,
  \item a \textbf{simple controller} that has become more forward-biased and less heading-driven,
  \item a \textbf{temporal localizer} that can now be reverted after jump rejection,
  \item a more faithful \textbf{MBRA inference wrapper},
  \item optional \textbf{depth veto / slowdown} support.
\end{itemize}
\end{mdframed}


\section{SDK Bridge: What It Does Now}

\subsection{Core idea}

The SDK side is now built as a \textbf{hybrid bridge}.

It no longer depends on only one path for command or telemetry flow.

\subsection{Control paths}

There are now two meaningful control paths:

\begin{enumerate}[nosep]
  \item \texttt{/control}
  \item \texttt{/control-legacy}
\end{enumerate}

\textbf{\texttt{/control}} uses the browser-backed path:
\begin{itemize}[nosep]
  \item the backend initializes a browser session,
  \item opens \texttt{/sdk},
  \item joins the SDK page,
  \item sends control through the JavaScript \texttt{window.sendMessage(...)} hook.
\end{itemize}

\textbf{\texttt{/control-legacy}} uses direct RTM send without requiring the
browser relay.

This means manual control and autonomous control are now more resilient to
browser-session failures than before.

\subsection{Telemetry paths}

Telemetry is also hybrid now.

The frontend JavaScript actively pushes data back into the backend:

\begin{itemize}[nosep]
  \item \texttt{basicRtm.js} listens to incoming RTM peer messages and posts them to \texttt{/api/update\_data}.
  \item \texttt{basicVideoCall.js} periodically captures the front camera frame and posts it to \texttt{/api/update\_frame}.
\end{itemize}

Then the backend exposes:

\begin{itemize}[nosep]
  \item \texttt{/data}
  \item \texttt{/v2/front}
\end{itemize}

Those endpoints now:

\begin{itemize}[nosep]
  \item prefer the lightweight cache populated by the JS push path,
  \item but fall back to \texttt{browser\_service} if the cache is empty.
\end{itemize}

\subsection{Practical consequence}

This is a major structural improvement over a pure pyppeteer-pull design.

The current SDK bridge is trying to achieve three things at once:

\begin{itemize}[nosep]
  \item low-latency data access through pushed caches,
  \item survivability when the lightweight cache is empty,
  \item a legacy fallback when the browser-backed command path is not healthy.
\end{itemize}

\subsection{Manual teleop}

The current intended manual teleop is \texttt{examples/keyboard\_control.py}.

Its behavior is:

\begin{itemize}[nosep]
  \item initialize \texttt{/sdk},
  \item capture live W/A/S/D key state with \texttt{pynput},
  \item smooth target velocities,
  \item send commands through \texttt{/control},
  \item fall back to \texttt{/control-legacy} if needed.
\end{itemize}

\texttt{examples/simple\_control.py} was also upgraded, but it is more of a
minimal terminal teleop helper than the main curses-based teleop path.


\section{Live Runtime: What It Does Now}

\subsection{Original role vs current role}

Originally, the live runtime was conceptually a thin loop:

\begin{center}
\fbox{\parbox{0.92\textwidth}{
camera $\rightarrow$ localize $\rightarrow$ plan $\rightarrow$ controller
$\rightarrow$ send command
}}
\end{center}

That is no longer the full story.

The current \texttt{live\_indoor\_runtime.py} contains a large amount of
\textbf{runtime policy} and \textbf{recovery logic} around the controller.

\subsection{What the live loop does now}

The loop still performs the same main phases:

\begin{enumerate}[nosep]
  \item get camera frame from the SDK,
  \item get telemetry from the SDK,
  \item run the motion-state filter,
  \item run localization,
  \item run graph planning,
  \item assemble controller input,
  \item ask the selected controller for a command,
  \item post-process that command through runtime safety / recovery logic,
  \item optionally send the command.
\end{enumerate}

But after controller inference, the runtime now adds several additional layers.

\subsection{Runtime layers added on top of the controller}

\begin{enumerate}[nosep]
  \item \textbf{Controller-specific defaults}:
  \item MBRA defaults to 3 Hz and 4 subgoal hops.
  \item Simple defaults to 2 Hz and 15 subgoal hops.
\end{enumerate}

\begin{enumerate}[resume,nosep]
  \item \textbf{Target reached confirmation}:
  \item target is not declared reached from one frame alone,
  \item it requires confidence and repeated confirmation ticks.
\end{enumerate}

\begin{enumerate}[resume,nosep]
  \item \textbf{Localization jump rejection}:
  \item large step jumps require higher confidence,
  \item rejected jumps trigger localizer state revert.
\end{enumerate}

\begin{enumerate}[resume,nosep]
  \item \textbf{Proximity slowdown}:
  \item forward speed is reduced near the target.
\end{enumerate}

\begin{enumerate}[resume,nosep]
  \item \textbf{Recovery backup mode}:
  \item a reverse override can temporarily replace normal control,
  \item after recovery, runtime and controller state can be reset.
\end{enumerate}

\begin{enumerate}[resume,nosep]
  \item \textbf{Angular saturation override}:
  \item repeated near-max angular output triggers a forced straight override.
\end{enumerate}

\begin{enumerate}[resume,nosep]
  \item \textbf{Optional monocular depth veto}:
  \item slows or stops forward motion if forward clearance looks small.
\end{enumerate}

\begin{enumerate}[resume,nosep]
  \item \textbf{RPM stall detection}:
  \item low RPM during commanded forward motion can trigger backup recovery.
\end{enumerate}

\begin{enumerate}[resume,nosep]
  \item \textbf{No-progress reset}:
  \item if the localized step does not change for long enough, the controller can be reset.
\end{enumerate}

\subsection{Most important conceptual consequence}

The controller file is no longer the whole control story.

The current on-robot behavior is a combination of:

\begin{itemize}[nosep]
  \item the selected controller,
  \item the sensor-state filter,
  \item localization confidence and jump rejection,
  \item runtime overrides,
  \item recovery logic,
  \item optional depth safety.
\end{itemize}

So when the robot behaves a certain way, it may not be coming from the
controller alone.


\section{Current Simple Controller: What It Really Is}

\subsection{What changed in spirit}

The current simple controller is much less like a continuous heading controller
than the original baseline intuition suggested.

It is now closer to:

\begin{center}
\fbox{\parbox{0.92\textwidth}{
brief alignment when needed $\rightarrow$ then mostly forward crawl
with outer runtime logic handling many failures
}}
\end{center}

\subsection{Important current settings}

Key changes in the current simple controller:

\begin{itemize}[nosep]
  \item \texttt{min\_linear} increased to 0.10,
  \item \texttt{drive\_heading\_gain} is now 0.0,
  \item alignment thresholds are much wider,
  \item alignment duration is capped by \texttt{max\_align\_ticks},
  \item no-heading mode is more aggressive and always maintains forward crawl.
\end{itemize}

\subsection{Very important runtime interaction}

The runtime now explicitly sets:

\begin{center}
\texttt{controller\_input["subgoal\_orientation"] = None}
\end{center}

That means the simple controller usually does \textbf{not} operate with real
subgoal heading information during the live loop.

In practice, this pushes it into the \texttt{no\_heading\_forward\_crawl} style
of behavior much more often.

\subsection{Implication}

The current simple path is best described as:

\begin{itemize}[nosep]
  \item localization says where we are,
  \item planning says which nearby graph step is ahead,
  \item the controller tries to move forward conservatively,
  \item the runtime wrapper provides many of the recovery behaviors.
\end{itemize}

This is not a rich trajectory-following controller at the moment.


\section{Motion-State Filter}

\subsection{Role}

\texttt{sensor\_state.py} is still a lightweight motion-prior layer, not a full
sensor-fusion backbone.

\subsection{Current signals used}

It currently derives:

\begin{itemize}[nosep]
  \item filtered heading from orientation,
  \item filtered heading-rate from gyro Z,
  \item filtered RPM mean from the latest RPM packet,
  \item telemetry freshness / stale status.
\end{itemize}

\subsection{Important change}

The stale logic no longer depends on ``did update() get called just now?''

It now tries to use the sensor timestamp to decide whether data is actually
fresh. That is a more correct design.

The timeout is also now more tolerant at 5 seconds.

\subsection{Practical effect}

This filter is not driving localization directly in a heavy fusion sense.
Instead, it provides:

\begin{itemize}[nosep]
  \item heading hints,
  \item turn-rate damping hints,
  \item RPM-based motion hints,
  \item a freshness stop signal.
\end{itemize}


\section{Localization and Temporal State}

\subsection{Corridor localizer}

The core localization backbone remains the same:

\begin{itemize}[nosep]
  \item preprocess frame,
  \item run CosPlace descriptor inference,
  \item retrieve nearest database matches,
  \item score candidates through the temporal localizer.
\end{itemize}

\subsection{New practical changes}

The corridor localizer now adds two deployment-focused capabilities:

\begin{itemize}[nosep]
  \item explicit device selection / CPU fallback for CUDA incompatibility,
  \item the ability to revert the last temporal update.
\end{itemize}

\subsection{Temporal localizer change}

The temporal localizer now supports:

\begin{itemize}[nosep]
  \item \texttt{save\_state()}
  \item \texttt{revert\_state()}
\end{itemize}

This is specifically there so runtime jump rejection can say:

\begin{quote}
``That large localization jump was not trusted, so roll back the temporal
state and keep the previous estimate.''
\end{quote}

\subsection{Meaning}

This is an important conceptual change.

Before, temporal localization was just a forward accumulator.
Now it is part of a runtime closed loop where:

\begin{itemize}[nosep]
  \item localization proposes a step,
  \item runtime judges whether that jump is believable,
  \item runtime can reject it and restore the previous temporal state.
\end{itemize}


\section{Navigation Runtime Handoff}

\texttt{navigation\_runtime.py} is still a fairly clean handoff layer.

Its main job is:

\begin{enumerate}[nosep]
  \item localize,
  \item plan path to target,
  \item load the subgoal image if available,
  \item package controller input.
\end{enumerate}

\subsection{Important portability fix}

It now repairs absolute image paths from graph artifacts when they were created
on another machine.

This is practical and important because many graph databases were generated on a
different system and carried over.


\section{Current MBRA Path}

\subsection{Main conceptual stance}

The code now treats MBRA as a \textbf{short-horizon local controller only}.

That matches the intended project architecture and also matches the strongest
handoff documents.

\subsection{What the current wrapper does}

The MBRA wrapper now:

\begin{itemize}[nosep]
  \item loads model config and weights,
  \item maintains only the observation image history,
  \item uses a fixed \texttt{vel\_past} tensor,
  \item uses timestep \texttt{[0,0]} as the immediate action,
  \item clamps output to the repo's safe command range,
  \item enforces a minimum forward command to avoid deadlock.
\end{itemize}

\subsection{Why this matters}

Earlier deployment-style hacks such as:

\begin{itemize}[nosep]
  \item feeding back previous predicted commands into \texttt{vel\_past},
  \item EMA blending over outputs,
  \item keeping controller-side motion history,
\end{itemize}

have been removed.

The current wrapper is therefore much closer to:

\begin{quote}
``run the MBRA model the way the original design intended, then only add small
deployment clamps and confidence gating.''
\end{quote}

\subsection{MBRA utility loading}

\texttt{mbra\_repo/deployment/utils\_logonav.py} was changed to lazy-load heavy
model families.

This is not a navigation-policy change. It is a deployment / dependency change.

Its practical value is:

\begin{itemize}[nosep]
  \item MBRA and IL paths can load with fewer unnecessary imports,
  \item unrelated heavyweight stacks do not need to be imported unless actually used.
\end{itemize}


\section{Depth Path}

\texttt{depth\_estimator.py} remains an optional helper around Depth Anything V2.

The most important practical change is not the inference logic itself, but the
checkpoint search:

\begin{itemize}[nosep]
  \item it now searches the actual HuggingFace-style metric indoor filenames,
  \item it also preserves some legacy filename variants.
\end{itemize}

In the live runtime, depth currently acts only as:

\begin{itemize}[nosep]
  \item a forward-clearance slowdown layer,
  \item a forward-clearance stop layer,
  \item a trigger for recovery backup.
\end{itemize}

It is not part of localization or planning.


\section{Most Important Practical Takeaways}

\begin{enumerate}[nosep]
  \item \textbf{The SDK path is much stronger than before.}
  Command and telemetry now both have redundancy.

  \item \textbf{The live runtime is now a major source of behavior.}
  Reading only the controller file is no longer enough to understand the robot.

  \item \textbf{The simple controller is currently a crawl-heavy baseline.}
  It is not doing rich heading-based local control in the live loop.

  \item \textbf{Localization remains the strongest validated subsystem.}
  The newer changes mainly make it more robust to deployment issues and jump rejection.

  \item \textbf{MBRA has been made more faithful to the original reference behavior.}
  That makes live evaluation more meaningful than before.

  \item \textbf{Many changes were practical deployment fixes, not architectural replacements.}
  The codebase has been hardened in-place rather than rebuilt from scratch.
\end{enumerate}


\section{Current Interpretation of the Stack}

\begin{mdframed}
\textbf{Best current reading of the codebase:}

\begin{itemize}[nosep]
  \item visual localization + temporal filtering + graph planning remain the backbone,
  \item the SDK bridge is now materially better than it was before,
  \item the runtime wrapper has become much more responsible for recovery and command shaping,
  \item the simple controller is conservative and forward-biased,
  \item MBRA is now implemented more faithfully as an optional short-horizon controller,
  \item the main uncertainty is not plumbing anymore, but real autonomous execution behavior.
\end{itemize}
\end{mdframed}


\section{Related Documents}
\begin{itemize}
  \item \texttt{erc3\_full\_documentation.tex} --- single master guide for the complete project story and current architecture.
  \item \texttt{live\_indoor\_runtime\_story.tex} --- indoor evolution, MBRA integration, and checkpoint-step runtime behavior.
  \item \texttt{live\_outdoor\_ultra\_marathon\_story.tex} --- outdoor and marathon runtime evolution with safety-layer reasoning.
  \item \texttt{outdoor\_perception\_review.tex} --- depth/semantic perception findings and their runtime implications.
\end{itemize}

\end{document}
